\section{Methodology}
\label{sec:methodology}

We describe our experimental setup for measuring steering vector decay dynamics,
including model selection, steering vector extraction, experimental conditions,
and evaluation metrics.

\subsection{Model and Infrastructure}
\label{sec:model}

We use \gptxl (1.5B parameters, 48 layers, $d_{\text{model}} = 1600$), loaded
via TransformerLens~\citep{nanda2022transformerlens} in float16 on an NVIDIA
RTX~A6000 GPU (49~GB VRAM). We select \gptxl because it is well-supported by
TransformerLens, fits on a single GPU, and was used in the original \actadd
paper~\citep{turner2023actadd}. All generation uses greedy decoding (argmax)
with a fixed random seed of 42 for reproducibility.

\subsection{Steering Vector Extraction}
\label{sec:extraction}

We construct contrastive pairs for three behavioral dimensions: truthfulness
(20~pairs of true vs.\ false factual statements), positive sentiment
(20~pairs of positive vs.\ negative evaluative language), and formal style
(20~pairs of formal vs.\ informal register). Each pair consists of semantically
similar sentences differing primarily in the target dimension (\eg ``The Earth
revolves around the Sun'' vs.\ ``The Sun revolves around the Earth'' for
truthfulness).

For each behavior and target layer $\ell$, we compute the mean activation of
the last token across all positive and negative examples, then take the
normalized difference:
\begin{equation}
\steeringvec^{(\ell)} = \frac{\bar{\vh}^{+}_{(\ell)} - \bar{\vh}^{-}_{(\ell)}}
{\|\bar{\vh}^{+}_{(\ell)} - \bar{\vh}^{-}_{(\ell)}\|}
\label{eq:sv_extraction}
\end{equation}
where $\bar{\vh}^{+}_{(\ell)}$ and $\bar{\vh}^{-}_{(\ell)}$ are the mean
activations at layer $\ell$ for the positive and negative examples, respectively.

\subsection{Experimental Conditions}
\label{sec:conditions}

We generate continuations from 25 diverse open-ended prompts (\eg ``The most
important thing to remember is that''), designed to be neutral with respect to
all behavioral dimensions. For each prompt, we run six conditions:

\begin{enumerate}[leftmargin=*,itemsep=1pt,topsep=2pt]
\item \textbf{No Steering} (baseline): Normal generation with no vector applied.
\item \textbf{Continuous Steering}: Vector $\alpha \cdot \steeringvec$ added at
      every token position, where $\alpha$ is the multiplier strength.
\item \textbf{AR} (Autoregressive): Vector applied for the first $K$ tokens,
      then free generation without the vector.
\item \textbf{TF} (Teacher-Forced, steered tokens): Vector applied for $K$
      tokens, then the same tokens that \AR generated are fed as input, but
      without the vector.
\item \textbf{TF-Clean} (Teacher-Forced, clean tokens): Vector applied for $K$
      tokens, then tokens from the unsteered baseline are fed without the vector.
\item \textbf{Random} (control): A random unit vector with the same multiplier
      applied for $K$ tokens, then free generation.
\end{enumerate}

The \TF condition tests whether the model retains any hidden state memory of
the steering intervention beyond what the generated tokens encode. The \TFclean
condition tests whether behavioral persistence requires the steered tokens to
remain in context.

\subsection{Metrics}
\label{sec:metrics}

At each token position $t$ after steering removal, we record four metrics:

\para{Cosine similarity delta ($\Delta$cos).}
We compute the cosine similarity between the hidden state $\vh_t$ at layer
$\ell$ and the steering vector $\steeringvec$, then subtract the corresponding
value from the unsteered baseline:
\begin{equation}
\Delta\cossim(t) = \cos(\vh^{\text{steered}}_t, \steeringvec) -
                    \cos(\vh^{\text{unsteered}}_t, \steeringvec)
\label{eq:delta_cos}
\end{equation}
This isolates the steering-specific effect from natural variation in hidden
state alignment.

\para{Projection magnitude.}
The scalar projection of the hidden state onto the steering direction:
$\text{proj}(t) = \vh_t \cdot \steeringvec$.

\para{KL divergence from unsteered baseline.}
$\KL(p_t^{\text{steered}} \,\|\, p_t^{\text{unsteered}})$, measuring how
different the output distribution at position $t$ is from the unsteered model.

\para{KL divergence from steered baseline.}
$\KL(p_t^{\text{condition}} \,\|\, p_t^{\text{continuous}})$, measuring
divergence from the continuously steered model.

\subsection{Default Hyperparameters}
\label{sec:hyperparams}

Unless otherwise noted, we use a steering multiplier of $\alpha = 3.0$,
injection at layer $\ell = 24$ (the middle layer), and steering duration
$K = 3$ tokens. These are standard values from the \caa
literature~\citep{panickssery2024caa}. We generate 20 post-steering tokens
for decay measurement.

\subsection{Ablation Design}
\label{sec:ablation_design}

We systematically vary three parameters:

\begin{itemize}[leftmargin=*,itemsep=1pt,topsep=2pt]
\item \textbf{Layer}: $\ell \in \{8, 16, 24, 32\}$, spanning early to late
      layers.
\item \textbf{Multiplier}: $\alpha \in \{1.0, 2.0, 3.0, 5.0\}$, covering
      weak to strong steering.
\item \textbf{Duration}: $K \in \{1, 3, 5, 10\}$, testing whether longer
      steering creates accumulated effects.
\end{itemize}

Each ablation varies one parameter while holding the others at their default
values, using the positive sentiment behavior and 10 prompts per condition.
